\documentclass[../main.tex]{subfiles}
\begin{document}
\setlist[itemize]{topsep=2pt,itemsep=1pt,parsep=0pt,partopsep=0pt}
\setlist[itemize,2]{topsep=1pt,itemsep=0pt,parsep=0pt,partopsep=0pt}
\setlength{\parskip}{3pt}

\chapter{Week 1 — Introduction to Research}
\label{ch:week1}

\textit{Snapshot and part summaries: see the overview on page~\pageref{ov:week1}.}

\section{Part 1 — The Research Process}
\label{sec:research-process}

\begin{itemize}
  \item Six steps: \keyterm{problem} $\to$ \keyterm{sample} $\to$ \keyterm{research design} $\to$ \keyterm{measure} $\to$ \keyterm{analyse} $\to$ \keyterm{conclude}.
  \item Problem is broad by design; sample is a necessary subset — representativeness and known limitations bound what the conclusion may claim.
  \item Design (what and why vs.\ how and when) is expanded in Week~2 (\autoref{sec:conceptual-technical}).
  \item Measurement takes four forms: experiments, analytical calculations, numerical simulations, questionnaires/focus groups — mapping onto data types in Week~3 (\autoref{sec:data-sources}).
  \item Analysis: data processing $\to$ grouping $\to$ interpretation. Raw observations are not yet evidence.
  \item Conclusion states results, compares to expectations, reports limitations, next steps, and relevance.
\end{itemize}

\section{Part 2 — The Language of Research}
\label{sec:language-research}

\begin{itemize}
  \item \keyterm{Hypothesis}: specific prediction of what the study expects to happen.
  \item \keyterm{Descriptive study}: documents the value of an unknown variable.
  \item \keyterm{Relational study}: investigates whether two or more variables are associated.
  \item \keyterm{Causal study}: determines whether one variable drives an outcome — demands a design that rules out third variables.
  \item \keyterm{Variable}: smallest research unit; its possible values are \keyterm{attributes} (exhaustive and mutually exclusive).
  \item Attributes are quantitative or qualitative; qualitative data can always be summarised numerically.
  \item Relationship pattern: none, positive, negative, or curvilinear. Nature: correlational, causal, or third-variable problem.
  \item \keyterm{Correlation is not causation}: association alone never establishes a causal arrow.
\end{itemize}

\section{Part 3 — The Hourglass: Deduction and Induction}
\label{sec:hourglass}

\begin{itemize}
  \item \keyterm{Hourglass model}: broad question $\to$ narrows to operationalisation $\to$ measure and analyse $\to$ widens to generalised conclusion.
  \item \keyterm{Deductive reasoning}: general $\to$ specific (down the funnel).
  \item \keyterm{Inductive reasoning}: specific observations $\to$ broader generalisation (up the funnel) — fallible; sample must support the generalisation (\autoref{sec:sampling}).
  \item Cautionary cases: inferring population proportions from one sample; inferring a universal law from two confirming instances.
\end{itemize}

\section{Part 4 — Where Research Ideas Come From}
\label{sec:research-ideas}

\begin{itemize}
  \item Sources: practical field problems, gaps in specific knowledge, and conclusions of prior studies.
  \item \keyterm{Feasibility}: trade-off between rigour and practicalities (duration, ethics, cooperation, cost, tool access).
  \item \keyterm{Literature review}: identifies related research and sets the project in its theoretical context — answers what is known, which theories/concepts are used, which methods exist, and which questions remain open.
  \item Search strategy: credible scientific literature, keywords for academic engines, mind map of authors and areas, follow references backward and citations forward from a strong baseline article (\autoref{sec:referencing}).
\end{itemize}

\end{document}
